% RESULTADO_RECUPERADO: funtor dimensional integral, eq:cZ-internos,
% eq:longitud-ruta y definición de ell_0; composición con artículo III.
% La prueba reunida y su control algebraico son específicos de esta contribución.
\subsection{Vitesse constitutive et réalisation cinématique de l’horloge}
\label{v:sec:enlace-velocidad-constitutiva}

La propagation radiative utilise la section de vitesse produite par le
transport angulaire du même état HMT. L’identification est établie
avant le choix des unités de représentation. Soient $x=A_{\mathrm{rad}}$ et
$y=C^*_{\mathrm{rad}}$ les coordonnées angulaires engendrées dans la
section~\ref{sec:ii-angulos}, avec les canaux
\eqref{eq:ii-canales}, et soit
$\mathsf R$ l’involution autoadjointe des deux feuillets, avec projecteurs
de rang un $P_+$ et $P_-$. Dans le domaine $x>|y|$, le transport et ses
canaux sont
\begin{equation}
 T=\exp(-xI-y\mathsf R)=q_+P_++q_-P_-,
 \qquad q_+=e^{-x-y},\quad q_-=e^{-x+y},\quad 0<q_\pm<1.
 \label{v:eq:velocidad-transporte}
\end{equation}
La réponse des blocs temporels de longueurs $90$ et $120$ est
\begin{equation}
 \mathsf V=(I-T^{90})(I-T^{120})^{-1}
       =r_+P_++r_-P_-,\qquad
 r_\pm=\frac{1-q_\pm^{90}}{1-q_\pm^{120}}.
 \label{v:eq:velocidad-respuesta}
\end{equation}
La construction part des canaux hérités ; les symboles $x,y$ et
$q_\pm$ désignent leurs coordonnées et non des paramètres ajustés à une vitesse
cible.

\begin{proposition}[Identité des réalisations constitutive et cinématique]
\label{v:prop:identidad-velocidad}
Dans un même système de coordonnées dimensionnelles, la section $c_{\mathrm{int}}$ de l’horloge
et la section $c_{\mathrm{vac}}$ de la réponse constitutive coïncident :
\begin{equation}
 c_{\mathrm{int}}=c_{\mathrm{vac}}
   =\widehat c\,u_C,
 \qquad
 \widehat c=(r_+r_-)^{-1}
            =\exp(-\mathcal E),
 \label{v:eq:identidad-velocidad}
\end{equation}
où
\begin{equation}
 \mathcal E=
 \log\!\bigl[(1-q_+^{90})(1-q_-^{90})\bigr]
 -\log\!\bigl[(1-q_+^{120})(1-q_-^{120})\bigr].
 \label{v:eq:velocidad-lector-simetrico}
\end{equation}
L’égalité conserve l’échange des feuillets.
\end{proposition}
\begin{proof}
Les facteurs qui apparaissent dans les logarithmes sont positifs. Par conséquent,
\[
 \mathcal E=\log(r_+r_-)
            =\log\det\mathsf V,
 \qquad \exp(-\mathcal E)=(\det\mathsf V)^{-1}.
\]
Le déterminant est calculé ici dans la réalisation à deux feuillets de rang
un. La définition cinématique $c_{\mathrm{int}}=\exp(-\mathcal E)u_C$
et la définition constitutive $c_{\mathrm{vac}}=(r_+r_-)^{-1}u_C$
sont, par conséquent, la même section. L’échange $P_+\leftrightarrow
P_-$ permute $r_+$ et $r_-$ et conserve leur produit. Enfin,
$\det T=e^{-2x}$ correspond au transport élémentaire, tandis que
$\det\mathsf V=r_+r_-$ correspond à sa réponse temporelle ; les deux
déterminants ont des fonctions distinctes dans cette composition.
\end{proof}

\subsubsection{Longueur de trajectoire et durée élémentaire}

Le repère dimensionnel positif $(u_S,u_C,u_T,u_Q)$ induit la base de longueur $u_L=u_Cu_T$. Pour une trajectoire enrichie $\gamma$ de $n=|\gamma|$ transitions et à canal constant, les lecteurs additifs sont
\begin{equation}
 T_{\mathrm{int}}(\gamma)=n u_T,
 \qquad L_{\mathrm{int}}(\gamma)=n\widehat c\,u_L.
 \label{v:eq:velocidad-longitud-reloj}
\end{equation}
Son additivité se rapporte au dénombrement et à la réalisation de chaque transition :
les chemins enrichis conservent en outre l’orientation et la mémoire.

\begin{corollary}[Normalisation de la longueur élémentaire]
\label{v:cor:longitud-elemental-constitutiva}
Pour $n>0$, $L_{\mathrm{int}}(\gamma)/T_{\mathrm{int}}(\gamma)
=c_{\mathrm{vac}}$. Dans la normalisation $t_0=u_T$,
\begin{equation}
 \ell_0=c_{\mathrm{int}}t_0=\widehat c\,u_L.
 \label{v:eq:longitud-elemental-constitutiva}
\end{equation}
Dans un système de coordonnées temporelles général $t_0=\tau_0u_T$, avec $\tau_0>0$,
l’expression correspondante est $\ell_0=\tau_0\widehat c\,u_L$.
\end{corollary}
\begin{proof}
La division de \eqref{v:eq:velocidad-longitud-reloj} simplifie l’entier
$n$ et utilise $u_L/u_T=u_C$. La formule élémentaire s’obtient
en multipliant la section de vitesse par $t_0$.
\end{proof}

La généralisation à un canal dépendant de l’arête conserve la formule
additive $L_{\mathrm{int}}(\gamma)=\sum_{e\in\gamma}\widehat c(e)u_L$.
Son quotient par $nu_T$ est la moyenne des vitesses de ces arêtes.
La réalisation radiative homogène emploie le canal constant déclaré
dans \eqref{v:eq:velocidad-longitud-reloj}.

\subsubsection{Changement de coordinatisation et normalisation adaptée}

\begin{proposition}[Covariance de l’identification]
\label{v:prop:covariancia-identidad-velocidad}
Soient $u_C'=d u_C$ et $u_T'=\eta u_T$, avec $d,\eta>0$.
La même vitesse et la même durée élémentaire ont pour coordonnées
\begin{equation}
 \widehat c'=d^{-1}\widehat c,
 \qquad \tau_0'=\eta^{-1}\tau_0,
 \qquad u_L'=d\eta u_L.
 \label{v:eq:velocidad-cambio-carta}
\end{equation}
En particulier, si $c_V=\widehat c_Vu_C^V$ et
$u_C^V=d u_C^{III}$, l’égalité des sections équivaut exactement à
$d\widehat c_V=\widehat c_{III}$.
\end{proposition}
\begin{proof}
On exprime chaque section dans les deux bases. On obtient
$\widehat c'u_C'=\widehat c u_C$ et
$\tau_0'u_T'=\tau_0u_T$. Leur produit satisfait
\[
 \tau_0'\widehat c' u_L'
 =\frac{\tau_0}{\eta}\frac{\widehat c}{d}
       d\eta u_L
 =\tau_0\widehat c u_L=\ell_0.
\]
La dernière équivalence résulte de la substitution de la relation entre les bases.
\end{proof}

Le choix adapté $d=\widehat c$ donne $u_C'=c_{\mathrm{vac}}$ et
$\widehat c'=1$. Ce choix exprime la section déjà construite dans une
base qui lui est adaptée. Le coefficient spectral dans le système de coordonnées interne reste
$(r_+r_-)^{-1}$ ; l’inversion spectrale utilise ce système de coordonnées interne.
Par exemple, en maintenant les bases d’action et de charge, les coordonnées
constitutives se transforment comme
\[
 \widehat\mu'=d\widehat\mu,
 \qquad \widehat\varepsilon'=d\widehat\varepsilon,
 \qquad
 \widehat\mu'\widehat\varepsilon'=(\widehat c')^{-2}.
\]
Dans le système de coordonnées adapté, $\widehat\mu'=r_-/r_+$ et
$\widehat\varepsilon'=r_+/r_-$, car les coordonnées internes sont
$\widehat\mu=r_-^2$ et $\widehat\varepsilon=r_+^2$.
